\title{Large Language Models Can Automatically Engineer Features for Few-Shot Tabular Learning}

\begin{abstract}
Large Language Models (LLMs) exhibit outstanding proficiency in solving complex and novel reasoning tasks, making them highly promising for tabular learning, which is essential in numerous real-world scenarios. This work introduces an innovative in-context learning framework called \textsf{FeatLLM}, designed to utilize LLMs as feature engineers. \textsf{FeatLLM} transforms raw input data into an enriched dataset tailored for superior tabular prediction. The newly generated features are then passed to a basic downstream predictive model, such as linear regression, enabling robust few-shot learning performance. Notably, at inference, \textsf{FeatLLM} solely leverages this simple predictive model with the identified features, foregoing the necessity for LLM invocation on individual samples. In contrast to existing LLM-based methodologies, \textsf{FeatLLM} obviates repeated LLM queries during inference, only requiring API access to the LLM and sidestepping prompt size restrictions. Across extensive tabular datasets from diverse fields, empirical results reveal that \textsf{FeatLLM} constructs superior rules, achieving a notable (10\% average) improvement over methods like TabLLM and STUNT.
\end{abstract}

\section{Introduction}
Large Language Models (LLMs) have recently emerged as powerful tools for producing apt responses to novel tasks without the necessity for specialized model fine-tuning~\cite{creswell2022selection,imani2023mathprompter,taylor2022galactica}. Their training on large-scale textual datasets endows them with deep prior knowledge that can serve as a substitute for domain expert input~\cite{Genomes2021,singhal2023large,wilcox2003role}. This has facilitated automation across diverse fields, including dialogue interpretation~\cite{finch2023leveraging} and program correction~\cite{fan2023automated}. By embedding a clear task description and a handful of exemplars in prompts, LLMs contextualize the problem and focus attention on salient features and pertinent instances~\cite{brown2020language,zhao2023survey}. Such capacities underscore the potential of LLMs to function as capable feature engineering agents.

Expanding on the application of LLMs’ prior knowledge and inference capabilities, their use in tabular learning has garnered interest. Domain insights, such as correlations between age and disease susceptibility, can substantially improve predictive tasks in medicine and beyond. Contemporary approaches encode task descriptions and feature semantics in natural language, structure the data, and supply it to an LLM~\cite{dinh2022lift,wang2023anypredict}. These processes are typically supplemented by in-context learning~\cite{nam2023semi} or by applying parameter-efficient tuning techniques~\cite{dinh2022lift,hegselmann2023tabllm}. The infusion of LLM-driven prior knowledge has demonstrated considerable advantages, consistently surpassing standard tabular learning methods in settings with limited data~\cite{hegselmann2023tabllm}.

Existing LLM-driven approaches for tabular data analysis face notable shortcomings. When performing end-to-end predictions, these methods necessitate at least one inference from the LLM for each individual sample, which introduces considerable computational overhead. Additionally, achieving satisfactory accuracy typically hinges on fine-tuning the LLM, a practice that becomes impractical if the training code or tuning APIs are unavailable—an increasing issue as state-of-the-art LLMs increasingly restrict access to inference-only APIs~\cite{anil2023palm,openai2023gpt4}. Moreover, the majority of these strategies struggle with handling extensive prompt lengths~\cite{liu2023lost}; as the number of attributes in tabular datasets increases, resulting in longer serialized inputs, their feasibility and performance deteriorate. Such limitations fundamentally obstruct the deployment of LLM-based tabular learning solutions in practical settings.

In contrast, our proposed method moves beyond the conventional use of LLMs for direct prediction. We leverage LLMs as tools for feature engineering, allowing them to apply their pre-existing knowledge more efficiently. We introduce a new in-context learning paradigm, \textsf{FeatLLM}, designed to elucidate the ``criteria" guiding LLM prediction processes. As an illustrative case, when predicting whether a patient has a specific disease, an LLM may articulate explicit conditions—rules connecting feature values to a diagnosis. These inferred criteria serve as new features, superseding the original inputs in downstream tasks. This strategy repositions the function of the LLM from prediction to feature generation, which simplifies the subsequent learning pipeline: after these features are formulated, the need for a complex model during inference is reduced, thereby decreasing latency. Utilizing LLMs' ability to learn from context—simply by including exemplar demonstrations within prompts—enables feature extraction without conventional model training. Additionally, adopting ensemble learning techniques~\cite{breiman2001random} such as feature bagging addresses the issue of oversized prompts, effectively circumventing input-length constraints.

\textsf{FeatLLM} organizes the input prompt to facilitate feature creation by splitting it into two distinct tasks: (1) grasping the core problem and discerning how features relate to targets; and (2) formulating discriminative rules that separate classes, utilizing both prior insights and illustrative few-shot samples. This stepwise, systematic methodology enables LLMs to disregard irrelevant features when generating rules. The derived rules are implemented on the dataset—treated as programmatic instructions—and subsequently, the dataset is encoded into binary features that indicate compliance with each rule for every instance. The resulting features are then used to train a simple model that estimates class probabilities. Feature generation is further enhanced in terms of robustness through ensemble methods, where repeated feature generation and bagging are employed. Adjusting the bagging parameters, such as the count of features and samples, helps circumvent limitations related to large prompt sizes. \textsf{FeatLLM} operates without the need for additional training and incurs minimal inference overhead, thus broadening the applicability of LLMs in practical scenarios.

We validate the efficacy of \textsf{FeatLLM} across 13 diverse tabular datasets in settings with limited sample availability, demonstrating its reliable and superior performance. Our experiments highlight how LLMs adeptly leverage both prior domain knowledge and the underlying information present in data to extract high-quality features. The proposed framework surpasses recent few-shot learning benchmarks in multiple scenarios. The implementation is publicly available at an anonymized GitHub repository:~\texttt{https://github.com/Sungwon-Han/FeatLLM}.

\section{Related Work}

\subsection{Few-Shot Learning with Tabular Data}
Few-shot learning methodologies have made it possible to derive robust, generalizable features from data when only limited labeled examples are available~\cite{chen2018closer}. While originally devised for image-based applications, these techniques have shown adaptability to other domains, such as text, audio, and most recently, tabular datasets~\cite{majumder2022few,nam2023stunt,schick2021exploiting}. A major challenge in few-shot scenarios is the heightened susceptibility to the model learning irrelevant or spurious associations~\cite{bartlett2021deep}. As a response, emerging research focuses on supplementing models with auxiliary information or embedding domain-relevant priors to create better inductive biases during training. For example, practitioners harness abundant unlabeled tabular datasets in combination with strategic augmentations as part of semi-supervised learning frameworks~\cite{ucar2021subtab,yoon2020vime}, or pursue pre-training schemes that are unsupervised and task-independent to obtain advantageous model initialization~\cite{somepalli2021saint}. Common unsupervised training objectives include techniques such as data masking and corruption followed by reconstructive learning~\cite{arik2021tabnet,majmundar2022met}, or deploying augmentation for contrastive learning~\cite{bahri2022scarf,chen2020simple}. Nevertheless, the intrinsic diversity within tabular datasets renders it difficult to design augmentation methods that are generally applicable, resulting in limited efficacy for these strategies in low-sample contexts~\cite{nam2023stunt}.

Recently, researchers have explored training models across a wide spectrum of tabular datasets, not restricted to those closely associated with the task of interest, before transferring them to new tasks~\cite{zhang2023generative,levin2022transfer,zhu2023xtab}. For example, TransTab~\cite{wang2022transtab} leverages transformer architectures to capture semantic column relationships, utilizing column names rather than just the raw tabular entries. Insights acquired from numerous tables and column contexts enable the model to carry out zero-shot and few-shot predictions for unseen tabular tasks. In contrast, TabPFN~\cite{hollmann2022tabpfn} pre-trains Bayesian neural networks using synthetic datasets generated from diverse distributions. Post-pre-training, predictions are produced by directly inputting the target task's training and testing samples, foregoing additional retraining. This accumulated prior knowledge from various tabular sources consistently outperforms classic tree-based methods and enhances performance in few-shot scenarios.

\subsection{Language-Interfaced Tabular Learning}
Integrating large language models (LLMs)~\cite{anil2023palm,openai2023gpt4} represents another avenue for leveraging external knowledge. LLMs, trained on broad and heterogeneous text collections, exhibit remarkable ability to extrapolate to new tasks, underscoring their cross-domain effectiveness~\cite{brown2020language}. Recent explorations have sought to exploit the comprehensive world knowledge stored in LLMs for tabular data analysis. This typically involves encoding tabular records as textual input, which is then provided to the LLM via prompts~\cite{dinh2022lift,hegselmann2023tabllm,wang2023anypredict}. With input prompts containing actual tabular records, explanatory feature descriptions, and task-specific instructions, models can infer interdependencies between features and tasks, facilitating predictive inference. Follow-up investigations either apply parameter-efficient fine-tuning techniques such as LoRA~\cite{hu2021lora} or IA3~\cite{liu2022few}, or opt for in-context learning where a handful of annotated examples are embedded in prompts without needing additional model training~\cite{dinh2022lift,hegselmann2023tabllm,nam2023semi,slack2023tablet}.

Although the techniques discussed above significantly enhance few-shot tabular learning, their dependence on routing inference through an LLM presents obstacles for practical adoption, particularly in industrial environments. The current work strives to move away from the typical black-box methodology where each sample’s prediction is handled in an end-to-end manner by the LLM. Rather than processing every individual sample via the LLM, this research proposes leveraging the LLM to deduce the specific conditions or decision rules tied to each class. \textsf{FeatLLM} converts these deduced rules into programmatic code, generating additional features. This approach allows for sample predictions without requiring the LLM at inference, and employs in-context learning to harness rules based on prior domain knowledge and a limited number of sample instances.

\section{Methods}
\textsf{FeatLLM} functions to derive “rules,” which are the conditions pertinent to each answer class, utilizing the provided few-shot samples as illustrated in Figure~\ref{fig:main_model}. These rules, extracted via the LLM, are then translated into binary features indicating, for each sample, the presence or absence of the rule. The resulting binary feature set is input to a dot product layer with non-negative trainable weights to yield an estimate for the probability that each class is correct. Through model training, one can determine the weight and relevance of each feature in influencing class probabilities (see Section~\ref{Sec:training}). The entire process is carried out several times using bagging, and results are aggregated through ensembling. Below, the formulation of the problem and specifics of each procedural step are detailed.

\vspace*{-0.12in}
\paragraph{Problem formulation.} We examine a tabular dataset containing $N$ labeled instances, expressed as $\mathcal{D} = \{(\mathbf{x}^i, \mathbf{y}^i)\}_{i=1}^{N}$. Each instance $\mathbf{x}^i$ is a $d$-dimensional feature vector, and each feature is identified by its natural language name within $F = \{f_j\}_{j=1}^{d}$; the corresponding feature definitions are listed separately. Within a classification framework, each label $\mathbf{y}^i$ belongs to one of several pre-established categories. In $k$-shot learning scenarios, we use a subset of $k$ ($<N$) samples drawn at random from the labeled data to train the model.

\subsection{Prompt Design for Extracting Rules}
\label{Sec:prompt_design}
To facilitate rule extraction by the LLM with enhanced reasoning fidelity, we structure the prompt to emulate the problem-solving approach of an expert engaged with a tabular dataset. Our prompt construction includes three primary elements, detailed in Fig.~\ref{fig:prompt_for_rule} and in Appendix~\ref{sec:example_rule}:

\vspace*{-0.12in}
\paragraph{Basic information description.} This section delivers the crucial information required to tackle the task (highlighted in orange in Fig.~\ref{fig:prompt_for_rule}). It encompasses a task description with label-related details, descriptions of each feature, and illustrated examples. The task description is framed as a question (for instance, ``Does the data concerning this patient's myocardial infarction complications suggest chronic heart failure? Yes or no?"). Further task description samples appear in Appendix~\ref{sec:task_desc}. Feature descriptions specify the value type, and their definition may also be included; for categorical features, possible category values can be listed. Demonstrative examples are generated by serializing several training samples into textual form, each paired with its true label. This serialization adopts formats from established works~\cite{dinh2022lift,hegselmann2023tabllm}:

\begin{align}
&\text{Serialize}(\mathbf{x}^i,\mathbf{y}^i, F) = \nonumber \\ 
&``\ f_1\ \text{is}\ \mathbf{x}^i_1.\ ... \ f_d\  \text{is}\  \mathbf{x}^i_d.\ \ \text{Answer:}\  \mathbf{y}^i\ ”
\end{align}

\vspace*{-0.12in}
\paragraph{Reasoning instruction.} Our approach seeks to improve the LLM’s reasoning capabilities by providing explicit reasoning instructions (refer to the blue annotations in Fig.~\ref{fig:prompt_for_rule}). The instructions commence with an introductory statement, mirroring the style of the chain-of-thought methodology~\cite{wei2022chain}. Subsequently, the inference process is split into two distinct phases. Initially, the LLM is tasked with discerning causal relationships or patterns between features and the task description, doing so independently of any demonstration examples, and instead drawing upon its general knowledge and commonsense understanding. This allows the model to fully leverage its pre-existing knowledge base regarding the problem. In the follow-up phase, the LLM synthesizes information from step one, along with demonstration examples, to derive rules that pertain to each class. This bifurcated reasoning framework aids in mitigating the detection of non-essential or misleading correlations in unrelated data columns and contributes to a concentration on more meaningful attributes. To avert issues arising from producing too many or too few rules—such as underfitting or unnecessarily lengthy prompts—a predetermined number of rules (specifically, 10)\footnote{A discussion on rule quantity is provided in Section~\ref{sec:analysis}.} is mandated per class. Moreover, when the LLM generates rules, it is prompted to reference the feature type from the basic information description section, ensuring that rule formulation avoids any mismatches of feature types.

\vspace*{-0.12in}
\paragraph{Response instruction.} In order to streamline the parsing and practical use of the generated rules, we further instruct the LLM on formatting its responses, using targeted response instructions (highlighted in yellow in Fig.~\ref{fig:prompt_for_rule}). These prompts specify both the necessary response format for each answer class (i.e., structure for the response) and the precise construction of each rule (i.e., structure for the rule), enabling the LLM to adopt the suitable format corresponding to the feature type in question. The absence of restrictive guidelines permits the model to combine rules utilizing logical connectors such as AND or OR. A sample output is demonstrated in Appendix~\ref{sec:outcome-rule}.

\subsection{Inferring Class Likelihood via Rules}
\label{Sec:training}

\paragraph{Rule-based feature extraction.} At this step, the rules formulated earlier are utilized to construct new binary features. For every class, these features indicate whether a sample meets any of the class-specific rules—yielding a '0' or '1' for each sample-rule pair. These indicators serve as predictors of the probability that a sample is associated with a particular class. Due to the natural language origins of the LLM-generated rules, an automated feature generation strategy requires a tailored parsing approach. Although rule generation was designed to enforce response formatting to facilitate parsing, the actual output from LLMs frequently contains minor syntactic inconsistencies~\cite{kovavc2023large}. For example, LLMs might use textual substitutes such as ``greater than" or ``less than" instead of the symbols ``$>$” and ``$<$,” complicating the parsing process.

To handle this syntactic noise effectively without engineering complex parsing logic, we capitalize on the LLM's own semantic comprehension and programming abilities. LLMs are capable of interpreting text semantically despite variances in syntax and are adept at generating code when prompted, owing to their code-centric training data. To exploit these strengths, our prompt incorporates a function definition, detailed descriptions of inputs and outputs, and the parsed rules, which is subsequently provided to the LLM. Example prompts and outputs can be seen in Figures~\ref{fig:prompt_for_parsing} and ~\ref{sec:example_parsing}-\ref{sec:example_parse_outcome} in the Appendix. The resulting code is executed within Python using the \texttt{exec()} function, enabling the automatic conversion of data into features.

\vspace*{-0.12in}
\paragraph{Estimating probability of class membership.} With binary rule-based features per class now available, a naive way to quantify a sample's likelihood of belonging to a particular class is by tallying the number of satisfied class-specific rules—this amounts to summing the respective binary indicators. Nevertheless, because some rules are more informative than others, it is essential to learn their relative influence from labeled training data. To achieve this, we employ a standard machine learning (ML) model that estimates the significance of each feature for each class. As an example, consider a linear model with no intercept term. Suppose for sample $\mathbf{x}^i$ and class $k$, the binary feature vector is represented as $\mathbf{z}_k^i \in \{0, 1\}^R$, with $R$ indicating the total rules for class $k$, and $c$ signifying the number of classes. The probability vector $\mathbf{p}^i$ for each class is computed via:

\begin{align}
\text{logit}_k^i &= \max(\mathbf{w}_k, 0) \cdot \mathbf{z}_k^i, \nonumber \\
\mathbf{p}^i &= \text{Softmax}({[\text{logit}_{1}^i, ..., \text{logit}_{c}^i]}). \label{eq:infer_prob}
\end{align}

In this context, $\mathbf{w}_k$ denotes the tunable parameters within the linear model, serving as indicators of each feature's relevance. By restricting these weights to the positive subspace, it is possible to guarantee that features produced by the LLM will not diminish a class’s probability during the learning phase. The resultant logit vector is subsequently transformed into class probabilities via the softmax function. The cross-entropy loss is minimized to adjust $\mathbf{w}_k$, leveraging a limited set of training examples.

\vspace*{-0.12in}
\paragraph{Ensembling with bagging.} To achieve the final prediction, the methodology involves multiple executions of the whole procedure, thereby constructing an ensemble of inference models. When $T$ trials yield trained weight sets $\{\mathbf{w}[t]\}_{t=1}^T$, the ensemble prediction aggregates class probabilities, $\mathbf{p}^i$, for each weight instance $\mathbf{w}[t]$ as described in Eq.~\ref{eq:infer_prob}, by taking their average.

To foster variation in rule creation for the ensemble, the LLM’s inference temperature is raised, inducing sufficient randomness. Furthermore, the ordering of few-shot exemplars is shuffled, taking advantage of the fact that altering this sequence can influence the LLM’s response~\cite{lu2022fantastically}\footnote{Further analysis regarding rule diversity is presented in Appendix~\ref{sec:diversity}}. Exploiting the diversity in outcomes leads to stronger prediction robustness, and bagging is applied to randomly select subsets of features or samples for each iteration; this approach becomes increasingly beneficial as the training dataset grows in size or complexity. The ensemble strategy confers two principal benefits. Firstly, misgenerated rules from certain LLM trials can be counteracted by correct rules from others, collectively improving framework stability—this phenomenon aligns with LLM self-consistency~\cite{wang2022self}, wherein rules frequently inferred across trials tend to be reliable. Secondly, the ensemble helps overcome constraints imposed by prompt length; for large tabular datasets, bagging enables manageable prompt sizes while maintaining model performance. Although a single rule set may fail to represent all feature relationships, the collective pool of weak learners assembled from numerous trials effectively offsets deficits from any incomplete feature or sample representation in individual trials.

\section{Experiments}
We assess \textsf{FeatLLM} across a diverse set of tabular datasets and carry out multiple analytical studies to explore the inner workings of our framework. For a more exhaustive examination, including experiments with different LLMs and qualitative assessments, please refer to the Appendix.

\subsection{Performance Evaluation}
\paragraph{Datasets.}
Our experimental setup features 13 datasets designed for either binary or multi-class classification. The datasets used are as follows: (1) Adult~\cite{asuncion2007uci}, used to forecast whether a person's annual income surpasses \$50,000; (2) Bank~\cite{moro2014data}, aimed at predicting customer subscription to a term deposit; (3) Blood~\cite{yeh2009knowledge}, for estimating if blood donors are likely to return; (4) Car~\cite{kadra2021well}, utilized for assessing car quality; (5) Communities~\cite{misc_communities_and_crime_183}, which estimates the crime rate within a particular region; (6) Credit-g~\cite{kadra2021well}, for classifying individuals as good or bad credit risks; (7) Diabetes\footnote{\texttt{kaggle.com/datasets/uciml/pima-indians-diabetes-database}}, predicting diabetes diagnoses; (8) Heart\footnote{\texttt{kaggle.com/datasets/fedesoriano/heart-failure-prediction}}, for predicting coronary artery disease; and (9) Myocardial~\cite{misc_myocardial_infarction_complications_579}, which classifies individuals based on chronic heart failure status.

Additionally, we select more contemporary and synthetic datasets that are less commonly discussed and were absent from LLM pre-training: (10) Cultivars, which evaluates grain yield predictions for soybean cultivars\footnote{\texttt{archive.ics.uci.edu/dataset/913}}; (11) NHANES, intended for age group classification using health records\footnote{\texttt{archive.ics.uci.edu/dataset/887}}; (12) Sequence-type, for recognizing four distinct sequence types—arithmetic, geometric, Fibonacci, or Collatz; (13) Solution-mix, which identifies whether the concentration in a four-solution mixture exceeds 0.5. The first two were contributed to the UCI Machine Learning Repository after September 2023, making it certain that the LLM API (gpt-3.5-turbo-0613) we employ had no access to them during pre-training. The last pair represent synthetic datasets, thus they eliminate the potential for memorization by the LLM API during evaluation.

Details on the scale and complexity for each dataset are found in Table~\ref{Tab:basic-desc}.
Each dataset includes straightforward attribute names and descriptions. Large-scale datasets such as `Communities' and `Myocardial' contain in excess of 100 attributes, presenting unique difficulties in light of the LLM's finite context window.

\vspace*{-0.12in}
\paragraph{Baselines.}
Our comparisons feature nine baseline methods. The initial three are traditional methods for tabular data analysis: (1) Logistic Regression (LogReg); (2) XGBoost~\cite{chen2016xgboost}; and (3) RandomForest~\cite{ho1995random}. The fourth and fifth baselines are designed for scenarios with access only to unlabeled data: (4) SCARF~\cite{bahri2022scarf}, and (5) STUNT~\cite{nam2023stunt}; it should be highlighted that, in practical contexts, acquiring substantial unlabeled datasets is often challenging. A sixth recent baseline, (6) TabPFN~\cite{hollmann2022tabpfn}, leverages synthetic datasets with varied distributions for pre-training. The remaining baselines utilize LLM-based methodologies: (7) In-context learning (In-context)~\cite{wei2022emergent}, (8) TABLET~\cite{slack2023tablet}, and (9) TabLLM~\cite{hegselmann2023tabllm}. In-context learning places several training samples directly within the input prompt without explicit parameter tuning. TABLET augments inference quality by injecting supplementary hints—rule-sets and prototypes—via an external classifier into the prompt. TabLLM applies parameter-efficient tuning approaches such as IA3~\cite{liu2022few}, training the LLM on few-shot data.

\vspace*{-0.12in}
\paragraph{Implementation details.}
\textsf{FeatLLM} adopts GPT-3.5 as the foundational LLM, although its framework allows for flexibility across different LLM architectures (refer to Appendix~\ref{sec:backbones} for evaluation on varied backbones). During LLM inference, we apply a temperature of 0.5 and set the top-p parameter to 1, corresponding to the default API configuration. The system uses 20 ensembles and extracts 10 distinct rules in its feature extraction process. The influence of these hyper-parameters is examined in Figure~\ref{fig:hyperparameter2} and elaborated further in Appendix~\ref{sec:hyper-parameter}. For training the linear model, we utilize the Adam optimizer with a learning rate of 0.01, running for 200 epochs. $k$-fold cross-validation is used to identify the optimal stopping point. When reproducing baseline methods, GPT-3.5 is employed for in-context learning baselines, while TabLLM leverages T0~\cite{sanh2022multitask} in alignment with its original setup. It is important to highlight that TabLLM confines LLM usage to models with open checkpoints, making GPT-3.5 unavailable. As for conventional machine learning baselines, including LogReg, XGBoost, and RandomForest, their hyper-parameters are determined through Grid-search alongside $k$-fold cross-validation. The value of $k$ is chosen as either 2 or 4 to ensure all classes are represented in the training set. Comprehensive implementation details are provided in Appendix~\ref{sec:baselines}.

\vspace*{-0.12in}
\paragraph{Results.}
Tables~\ref{tab:main1} and \ref{tab:main2} detail comparative results over a diverse set of datasets. Each experiment was conducted three times with random splits of training and test data, reporting mean and standard deviation for the AUC metrics. In nearly all settings, \textsf{FeatLLM} either matches or surpasses other baseline models, typically occupying the leading or second rank in performance. Notably, our method attains results on par with traditional tabular baselines using just 4 shots, as opposed to the 32 shots these baselines require (refer to Fig.~\ref{fig:compares}). Despite instances where STUNT and TabPFN achieved notable gains on individual datasets, their aggregate advantage over standard machine learning approaches was limited. Furthermore, LLM-driven methods—including in-context learning, TABLET, and TabLLM—could not effectively infer on tabular datasets with a large number of features. By contrast, our framework leverages both bagging and ensembling, remaining effective and yielding strong outcomes even on high-dimensional datasets. Crucially, the core ideas behind our bagging and ensemble strategy could theoretically extend to any LLM-based technique. Nonetheless, such integration would necessitate twice as many inference steps per sample, dramatically increasing computation demands and, consequently, proving to be computationally infeasible in practice.

\subsection{Analysis \& Discussion}
\label{sec:analysis}
\paragraph{Ablation study.} We investigate the individual effect of primary framework components through a series of ablation experiments: (1) \textsf{-Tuning}—removing the parameter optimization step in the linear model and replacing it by summing binary features for logit computation; (2) \textsf{-Ensemble}—excluding the ensemble module; (3) \textsf{-Description}—eliminating feature descriptions; and (4) \textsf{-Reasoning}—removing the Step-1 reasoning procedure in rule generation.

Table~\ref{tab:ablation} displays average shifts in AUC scores after each ablation across all datasets considered. Performance consistently declined with the removal or modification of each constituent part. Remarkably, the positive impact of weight tuning was more pronounced at higher shot numbers, signifying that greater data volumes permit more precise rule weighting, thereby magnifying tuning’s contribution. In contrast, feature description and reasoning instruction played larger roles when fewer shots were available—pointing to the critical value of leveraging LLM's embedded prior knowledge in few-shot regimes. Finally, the proposed ensemble approach produced robust performance gains across all experimental conditions.

\vspace*{-0.12in}
\paragraph{Training \& inference time comparison.} We examine and contrast the training and inference durations across various approaches. All experiments utilize the Adult dataset, which includes a training set consisting of 16 instances. An A100 GPU serves as the standard computing resource, except in the case of TabLLM, which operates across four A100 GPUs to enable model parallelism. Table~\ref{tab:cost} details the runtime performance across methods. Importantly, \textsf{FeatLLM} achieves inference speeds that are on par with traditional approaches, indicating its efficiency. By comparison, alternative LLM-based techniques—including In-context learning, TABLET, and TabLLM—demand substantially greater inference time. With regard to training duration, \textsf{FeatLLM} aligns with the training times of other LLM-based approaches, necessitating merely 30 API queries. This relatively small number does not lead to considerable financial cost and is amenable to parallel execution. 

\vspace*{-0.12in}
\paragraph{Quality of features generated by \textsf{FeatLLM}.}
To assess the practical utility of rule-based features created by \textsf{FeatLLM}, we perform a basic experiment that evaluates their informativeness. In this setting, we determine the average AUROC score between each generated feature and the target labels, subsequently calculating the proportion of features with AUROC greater than 0.5 in every dataset. 
A summary of these findings is presented in Table~\ref{tab:quality}. The data demonstrates that most constructed rules pertain meaningfully to the target attribute, offering valuable information for tackling the associated task (see the ``Before tuning" column). Furthermore, when a linear model is tuned using available data, features with minimal relevance (specifically, those whose learned weights fall below a chosen threshold) are systematically excluded (see the ``After tuning" column). The threshold for exclusion was established at 0.1, based on the overall distribution of learned weights.

\vspace*{-0.12in}
\paragraph{Impact of introducing spurious correlations.} To evaluate the ability of different approaches to handle extraneous spurious correlations under conditions of limited labeled data, we carried out a dedicated assessment. In this experiment, we took the Heart dataset and randomly appended columns sourced from the Adult dataset which had no bearing on the Heart dataset’s objective. The subsequent variation in model performance was tracked as we incrementally increased the count of irrelevant columns. For consistency, we fixed the number of labeled samples at 8 and excluded scenarios with access to unlabeled data, thereby omitting methods like SCARF and STUNT from our evaluation.

Figure~\ref{fig:spurious} presents the impact of these artificially induced correlations on model performance. Among the methods, \textsf{FeatLLM} experienced the smallest reduction in accuracy resulting from the spurious correlations. This likely stems from our approach, which incorporates prior knowledge when deriving rules, explicitly linking features to the target task and discouraging the formulation of rules based on extraneous columns—resulting in greater robustness. Notably, rules relying on the noise-adding columns comprised a mere 1-2\% of all extracted rules. Nevertheless, \textsf{FeatLLM} is not immune to mistakenly learning spurious dependencies or misinterpreting causal links, especially when columns from semantically similar datasets are added at random (see Appendix~\ref{sec:spurious-appendix}).

\vspace*{-0.12in}
\paragraph{Hyper-parameter analysis}
\textsf{FeatLLM} incorporates two central hyper-parameters: the ensemble count and the number of rules per class generated from each LLM query. To assess their influence on model performance, we performed systematic experiments. Our findings indicate that increasing the ensemble size leads to better performance, but these improvements plateau after approximately 20 ensembles (refer to Fig.~\ref{fig:appendix-hyperparameter1} in Appendix). Regarding the rule count, although no significant statistical difference was observed, choosing 10 rules achieves the best balance between prompt size and model accuracy (see Fig.~\ref{fig:hyperparameter2}). We hypothesize this is because additional rules expand the input dimensions and thus the informational content, but in low-shot scenarios, this may promote overfitting.

\section{Conclusion}
This paper introduces \textsf{FeatLLM}, which establishes a new benchmark in few-shot tabular learning using LLMs. Distinct from prior approaches that rely on LLMs solely for direct predictions per sample, \textsf{FeatLLM} leverages LLMs to synthesize predictive criteria, adopting a feature engineering perspective. By integrating rule creation with an ensemble bagging methodology, \textsf{FeatLLM} minimizes inference costs, functions via simple API access without training, and bypasses restrictions associated with data feature dimensions. In summary, \textsf{FeatLLM} not only delivers strong predictive capabilities but also presents a highly effective, data-efficient method for deploying LLMs in practical tabular data settings.

\textsf{FeatLLM} currently targets exclusively the low-shot learning context. Moving forward, we aim to broaden its scope by enabling feature generation for datasets containing more samples, diversifying feature extraction strategies beyond rule-based methods, and enhancing interpretability, thereby facilitating application to a wider range of domains and tasks.

\section{Broader Impact}
The framework put forth—\textsf{FeatLLM}—is intended to harness LLMs’ prior knowledge for improved tabular learning outcomes, surpassing traditional supervised approaches. In advancing data-efficient learning, our methodology offers valuable solutions to overfitting problems arising from limited annotated data by supplementing models with external prior knowledge. \textsf{FeatLLM} holds particular promise for industry sectors such as finance and healthcare, where annotation is either costly or otherwise difficult, and pretrained LLMs encapsulate relevant domain information. An important avenue for future work entails systematic evaluations of how societal biases and inaccuracies within LLM prior knowledge influence downstream model predictions, as this prior knowledge may significantly affect, or even outweigh, the impact of the available labeled data.

\end{document}